\documentclass[a4j, 9pt]{ltjsarticle}

\usepackage{../mypackage}

\begin{document}

\section{対称式}
  対称式とは変数を入れ替えても形が変わらない多項式のことである。まず多項式がどんなものか確認して、次に対称式の性質や対称式の基本定理について議論しよう。  

  \subsection{多項式・整式の定義}
    \textbf{多項式(整式)}とは「数」と「変数」を和と積によって作られる
      (線形結合\footnote{
        線形結合の具体的な定義は
        「体$\ds \mathbb{F}$上の線形空間\boldsymbol{V}に対する$\ds \boldsymbol{v_1}, \boldsymbol{v_2}, \cdots , \boldsymbol{v_n} \in \boldsymbol{V}$の線形結合とは
        $\ds c_1 \boldsymbol{v_1} + c_2 \boldsymbol{v_2} + \cdots + c_n \boldsymbol{v_n} (c_1, c_2, \cdots , c_n \in \mathbb{F})$で表されるベクトルのこと」
      }と言ったりする)
    式のことである。\par
    たとえば$\ds 3x^3 + 5x^2 - x + 10$は$\ds x$を変数とする多項式である。
    また、多項式は変数を複数持つ場合も言う。\par
    変数が$\ds 1$つの場合の多項式の一般形を書くと

    \begin{equation*}
      \begin{split}
        $\ds a_n x^n + a_{n-1} x^{n-1} + \cdots + a_1 x + a_0 ( = \sum_{k=0}^n a_k x^k )$\\
      \end{split}
    \end{equation*}

    ここに各々の$\ds x_k$において$\ds a_k x^k$のようにして積で結ばれる。この積で結ばれた塊を\textbf{項}と呼び、$\ds a_k$を\textbf{$\ds x^k$の係数}と呼ぶ。
    特に$\ds 0$次の項を\textbf{定数項}と呼ぶ。\par
    また、一般には和を持たない多項式(例えば$\ds 4x^2$など)も存在し、それらを\textbf{単項式}と呼ぶ。\par
    ゆえに$\ds 単項式 \subset 多項式 = 整式$である。

    \subsubsection{多項式・整式でないもの}
      紛らわしいものに次のようなものがある。

      \begin{equation*}
        \begin{split}
          &・$\ds \frac{1}{x} = x^{-1}$\\
          &・$\ds x^\frac{1}{2} = \sqrt{x}$\\
          &・$\ds \sum_{n=0}^\infty x^n = 1 + x + x^2 + \cdots$\\
        \end{split}
      \end{equation*}

      これらはすべて多項式・整式ではない。(ゆえに単項式でもない)\par
      これらは上から順に「有利式・分数式」、(名前なし)、「級数」という。
  
  \columnbreak

  \subsection{対称式の定義}
    先に述べたように対称式とは変数を入れ替えても変わらない多項式のことである。
    例えば次のような$\ds 2$変数の多項式のようなものをいう。

    % \begin{equation}
    \begin{equation*}
      \begin{split}
        % $\ds f(x, y) \ldef x^2 + xy + y^2$ \label{sym_1}\\
        $\ds f(x, y) \ldef x^2 + xy + y^2$\\
      \end{split}
    % \end{equation}
    \end{equation*}

    この式で$\ds x, y$を入れ替えてみても次の通り

    % \begin{equation}
    \begin{equation*}
      \begin{split}
        % $\ds f(y, x) = y^2 + yx + x^2$ \label{sym_2}\\
        $\ds f(y, x) $  &$\ds = y^2 + yx + x^2$\\
                        &$\ds = f(x, y)$\\
      \end{split}
    % \end{equation}
    \end{equation*}

    % \eqref{sym_1}と\eqref{sym_2}は全く同じ多項式である。
    これらは全く同じ多項式である。\par
    今見たものは$\ds 2$変数のものであったが、さらには例として次のような$\ds n$変数の対称式を考えられる。

    \begin{equation*}
      \begin{split}
        $\ds f(x_1, x_2, \cdots , x_n) \ldef \sum_{k=1}^n x_k^7$\\
      \end{split}
    \end{equation*}

    これも任意の$\ds x_k (1 \leq k \leq n)$を任意の分だけ入れ変えてもすべての場合で同じ多項式となる。\par
    例えば$\ds x_2$と$\ds x_n$を入れ替えてみても

    \begin{equation*}
      \begin{split}
        $\ds f(x_1, x_n, \cdots , x_2) $  & $\ds = x_1^7 + x_n^7 + \sum_{k=3}^{n-1} x_k^7 + x_2^7$\\
                                          & $\ds = x_1^7 + x_2^7 + \sum_{k=3}^{n-1} x_k^7 + x_n^7$\\
                                          & $\ds = \sum_{k=1}^{n} x_k^7$\\
                                          & $\ds = f(x_1, x_2, \cdots , x_n)$
      \end{split}
    \end{equation*}

    のように成立する。

  \columnbreak

  \subsection{基本対称式の定義}
    基本対称式とは$\ds n$個の変数$\ds x_1, x_2, \cdots , x_n$らを要素にもつ集合$\ds X$、つまり$\ds X \ldef \{ x_k \mid k \in \mathbb{N} \land 1 \leq k \leq n \}$なる集合$\ds X$から
    $\ds 1 \leq k \leq n$を満たす異なる$\ds k$個の変数を掛け合わせたものをすべての場合で作り、それらの和をとってできる対称式のことをいう。
    この基本対称式は後に議論する「対称式の基本定理」に登場し、対称式に関してとても重要な役割を果たすため天から降ってきたようなやり方だがここでとりあえず定義しておく。\par
    この$\ds k$個の変数の場合の基本対称式を$\ds \sigma_k(x_1, x_2, \cdots , x_n)$とすると以下すべてが該当する。

    \begin{equation*}
      \begin{cases}
        &\begin{split}
          $\ds \sigma_1(x_1, x_2, \cdots , x_n)$ & $\ds = x_1 + x_2 + \cdots + x_n$\\
          &$\ds = \sum_{k=1}^n x_k$
        \end{split}\\
        &\begin{split}
            $\ds \sigma_2(x_1, x_2, \cdots , x_n)$ & $\ds = x_1 x_2 + x_1 x_3 + \cdots + x_1 x_n + x_2 x_3 + x_2 x_4 + \cdots + x_{n-1} x_n$\\
          & $\ds = \sum_{k_1=1}^{n-1} \sum_{k_2=2}^{n} x_{k_1} x_{k_2}$
        \end{split}\\
        &\begin{split}
          $\ds \sigma_3(x_1, x_2, \cdots , x_n)$ & $\ds = x_1 x_2 x_3 + x_1 x_2 x_4 + \cdots + x_{n-2} x_{n-1} x_n$\\
          & $\ds = \sum_{k_1=1}^{n-2} \sum_{k_2=2}^{n-1} \sum_{k_3=3}^{n} x_{k_1} x_{k_2} x_{k_3}$\\
          & $\ds = \sum_{k_1=1}^{n-2} \sum_{k_2=2}^{n-1} \sum_{k_3=3}^{n} \prod_{m=1}^3 x_{k_m}$
        \end{split}\\
        & \vdots \\
        &\begin{split}
          $\ds \sigma_l(x_1, x_2, \cdots , x_n)$ & $\ds = \sum_{k_1=1}^{n-l+1} \cdots \sum_{k_{l-1}=l-1}^{n-1} \sum_{k_l=l}^n \prod_{m=1}^l x_{k_m}$
        \end{split}\\
        & \vdots \\
        &\begin{split}
          $\ds \sigma_n(x_1, x_2, \cdots , x_n)$ & $\ds = x_1 x_2 x_3 \cdots x_n$\\
          & $\ds = \prod_{k=1}^n x_k$
        \end{split}\\
      \end{cases}
    \end{equation*}

    \subsubsection{基本対称式の具体例}
      これだけ見てもピンとこないだろうから簡単に$\ds n=2, 3$の場合を見てみよう。

      \paragraph{具体例$\ds 1$}
        $\ds n=2$の場合対称式の説明をしたときと照らせ合わせるように$\ds x_1 = x, x_2 = y$とすると

        \begin{equation*}
          \begin{cases}
            &$\ds \sigma_1(x, y) = x + y$\\
            &$\ds \sigma_2(x, y) = xy$\\
          \end{cases}
        \end{equation*}

        確かにこれらは対称式である。

        \subparagraph{類似: $\ds 2$次方程式の解と係数の関係}

          \begin{equation*}
            \begin{split}
              $\ds 2$次方程式$\ds ax^2 + bx + c = 0 (\Longleftrightarrow x^2 + \frac{b}{a}x + \frac{c}{a} = 0)$\\
            \end{split}
          \end{equation*}
          
          の$\ds 2$解を$\ds x= \alpha, \beta$としたとき、解と係数の関係として

          \begin{equation*}
            \begin{split}
              $\ds \alpha + \beta = - \frac{b}{a} \land \alpha \beta = \frac{c}{a} \Longleftrightarrow x^2 + \frac{b}{a}x + \frac{c}{a} = 0$\\
            \end{split}
          \end{equation*}

          が成立する。

      \paragraph{具体例$\ds 2$}
        次に$\ds n=3$の場合も同様に$\ds x_1 = x, x_2 = y, x_3 = z$として

        \begin{equation*}
          \begin{cases}
            &$\ds \sigma_1(x, y, z) = x + y + z$\\
            &$\ds \sigma_2(x, y, z) = xy + yz + zx$\\
            &$\ds \sigma_3(x, y, z) = xyz$\\
          \end{cases}
        \end{equation*}

        これも対称式であることが確認できた。

        \subparagraph{類似: $\ds 3$次方程式の解と係数の関係}

          \begin{equation*}
            \begin{split}
              $\ds 3$次方程式$\ds ax^3 + bx^2 + cx + d = 0 (\Longleftrightarrow x^3 + \frac{b}{a}x^2 + \frac{c}{a}x + \frac{d}{a} = 0)$\\
            \end{split}
          \end{equation*}

          の$\ds 3$解を$\ds x= \alpha, \beta, \gamma$としたとき、解と係数の関係として

          \begin{equation*}
            \begin{split}
              &$\ds \alpha + \beta + \gamma = - \frac{b}{a} \land \alpha \beta + \beta \gamma + \gamma \alpha = \frac{c}{a} \land \alpha \beta \gamma = \frac{d}{a}$\\
              &$\ds \quad \Longleftrightarrow x^3 + \frac{b}{a}x^2 + \frac{c}{a}x + \frac{d}{a} = 0$\\
            \end{split}
          \end{equation*}

          が成立する。

  \subsection{類似: 解と係数の関係と対称式}
    具体例$\ds 1,2$の類似で見たように、$\ds 2$次方程式$\ds ax^2 + bx + c = 0$の$\ds 2$解を$\ds x = \alpha, \beta$としたとき

    \begin{equation*}
      \begin{cases}
        &$\ds \sigma_1(\alpha, \beta) = -\frac{b}{a}$\\
        &$\ds \sigma_2(\alpha, \beta) = \frac{c}{a}$\\
      \end{cases}
    \end{equation*}

    同様に$\ds 3$次方程式$\ds ax^3 + bx^2 + cx + d = 0$の$\ds 3$解を$\ds x = \alpha, \beta, \gamma$としたとき

    \begin{equation*}
      \begin{cases}
        &$\ds \sigma_1(\alpha, \beta, \gamma) = -\frac{b}{a}$\\
        &$\ds \sigma_2(\alpha, \beta, \gamma) = \frac{c}{a}$\\
        &$\ds \sigma_3(\alpha, \beta, \gamma) = -\frac{d}{a}$\\
      \end{cases}
    \end{equation*}

    これを一般化すれば$\ds n ( \in \mathbb{N}$ ) 次方程式$\ds \sum_{k=0}^n a_k x^k = 0$の$\ds n$解を$\ds x = x_1, x_2, \cdots , x_n$とするとき

    \begin{eqnarray}
      \begin{cases}
        &$\ds \sigma_1(x_1, x_2, \cdots , x_n) = -\frac{a_{n-1}}{a_n}$\\
        &$\ds \sigma_2(x_1, x_2, \cdots , x_n) = \frac{a_{n-2}}{a_n}$\\
        &\vdots \\
        &$\ds \sigma_k(x_1, x_2, \cdots , x_n) = (-1)^k \frac{a_{n-k}}{a_n}$\\
        &\vdots \\
        &$\ds \sigma_n(x_1, x_2, \cdots , x_n) = (-1)^n \frac{a_0}{a_n}$\\
      \end{cases} \label{eq:solution-and-coefficients}
    \end{eqnarray}

    が言える。

  \subsection{多項式の辞書式順序の定義}
    単項式$\ds X, Y$を

    \begin{equation*}
      \begin{cases}
        &$\ds X \ldef A x_1^{p_1} x_2^{p_2} \cdots x_n^{p_n} = A \prod_{k=1}^n x_k^{p_k}$\\
        &$\ds Y \ldef B x_1^{q_1} x_2^{q_2} \cdots x_n^{q_n} = B \prod_{k=1}^n x_k^{q_k}$\\
      \end{cases}
    \end{equation*}

    と定義する。このとき次のような並べ方を\textbf{辞書式順序}と定義する。

    \begin{breakbox}
      \begin{equation*}
        \begin{split}
          &$\ds X$は$\ds Y$より強い($\ds Y$は$\ds X$より弱い)\\
          &\defineProposition
          \begin{cases}
            &$\ds p_1 > q_1$\\
            &$\ds p_1 = q_1 \land p_2 > q_2$\\
            &$\ds p_1 = q_1 \land p_2 = q_2 \land p_3 > q_3$\\
            &\vdots \\
            &$\ds p_1 = q_1 \land p_2 = q_2 \land$\\
            & \quad $\ds \cdots \land p_{n-1} = q_{n-1} \land p_n > q_n$\\
          \end{cases}\\
        \end{split}
      \end{equation*}
    \end{breakbox}

    \footnote{
      ここに用いた
      \begin{cases}
        &$\ds P_1(x)$\\
        &$\ds P_2(x)$\\
      \end{cases}は$\ds P_1(x) \lor P_2(x)$の意味
    }

    具体的には、多項式$\ds f(x_1, x_2, \cdots , x_n), g(x_1, x_2, \cdots , x_n)$において
    
    \begin{equation*}
      $\ds f$が$\ds g$より強い \quad $\ds \Longleftrightarrow$
      \begin{matrix}
        &$\ds f$の項のうち最も強い項が\\
        & $\ds g$の項のうち最も強い項より強い
      \end{matrix}
    \end{equation*}
    
  \subsection{対称式の基本定理}
    対称式には次のような定理がある。

    \begin{breakbox}
      \begin{equation*}
        \begin{split}
          &$\ds x_1, x_2, \cdots , x_n$についての対称式$\ds f(x_1, x_2, \cdots , x_n)$は\\
          &基本対称式$\ds \sigma_1, \sigma2, \cdots, \sigma_n$に関する多項式\\
          &$\ds g(\sigma_1, \sigma_2, \cdots , \sigma_n)$として一意に表せる。
        \end{split}
      \end{equation*}
    \end{breakbox}

    \subsubsection{対称式の基本定理を利用した具体例}
      例えば$\ds x, y$に関する対称式$\ds x^2 + y^2$は

      \begin{equation*}
        \begin{split}
          $\ds x^2 + y^2 $  &$\ds = (x + y)^2 - 2xy$\\
                            &$\ds = \sigma_{1}^2 - 2\sigma_2$\\
        \end{split}
      \end{equation*}

      また、$\ds \frac{y}{x} + \frac{x}{y}$の場合も

      \begin{equation*}
        \begin{split}
          $\ds \frac{y}{x} + \frac{x}{y}$ &$\ds = \frac{x^2 + y^2}{xy}$\\
                                          &$\ds = \frac{(x + y)^2 - 2xy}{xy}$\\
                                          &$\ds = \frac{\sigma_{1}^2 - 2\sigma_2}{\sigma_2}$
        \end{split}
      \end{equation*}

      と基本対称式で一意に表せた。

    \begin{proof}[\textbf{対称式の基本定理の証明}]
      任意の対称式が基本対称式で一意に表せるのはなぜか。証明してみよう。\par
      対称式$\ds f(x_1, x_2, \cdots , x_n)$に関して先ほど定義した辞書式順序的に最も強い項を

      \begin{eqnarray}
        \begin{split}
          $\ds A x_1^{p_1} x_2^{p_2} \cdots x_n^{p_n}$ \tag{$*$} \label{eq:strongest-paragraph}
        \end{split}
      \end{eqnarray}

      とおく。辞書式順序の強弱の定義より

      \begin{equation*}
        \begin{split}
          $\ds p_1 \geq p_2 \geq \cdots \geq p_n$
        \end{split}
      \end{equation*}

      である。また、$\ds n$個の$\ds x$に関する基本対称式を$\ds \sigma_1, \sigma_2, \cdots , \sigma_n$として\par

      \begin{equation*}
        \begin{split}
          $\ds g_1(\sigma_1, \sigma_2, \cdots , \sigma_n)$  &$\ds \ldef A \sigma_{1}^{p_1-p_2} \sigma_{2}^{p_2-p_3} \cdots \sigma_{n}^{p_n}$\\
                                                            &$\ds = A \left( \sum_{i=1}^n x_i \right)^{p_1-p_2} \left( \sum_{i=1}^{n-1} \sum_{j=2}^n x_i x_j \right)^{p_2 - p_3} \cdots (x_1 x_2 \cdots x_n)^{p_n}$
        \end{split}
      \end{equation*}

      のようにして基本対称式より作られる対称式$\ds g_1$を定義する。\par
      今、$\ds f$と$\ds g_1$を独立に定義した。ただ、どちらも$\ds x_1, x_2, \cdots , x_n$に関する対称式として定義しているので
      両者とも最も強い項は\eqref{eq:strongest-paragraph}と表される項である。\par
      ゆえに

      \begin{equation}
        \begin{split}
          $\ds f_1 \ldef f(x_1, x_2, \cdots , x_n) - $&$\ds g_1(\sigma_1, \sigma_2, \cdots , \sigma_n)$ \label{eq:define-formula-of-f_1}
        \end{split}
      \end{equation}

      とすれば$\ds f$より弱い対称式$\ds f_1$が定義できる。\par
      次に$\ds f_1$の最強の項について同様の操作をする。先ほど$\ds f$に対する対称式$\ds g_1$を定義したときと同様に今度は$\ds f_1$に対する対称式$\ds g_2$を定義して

      \begin{equation*}
        \begin{split}
          $\ds f_2 \ldef f_1(x_1, x_2, \cdots , x_n) - g_2(\sigma_1, \sigma_2, \cdots , \sigma_n)$
        \end{split}
      \end{equation*}

      のようにすることで$\ds f_1$より弱い対称式$\ds f_2$を定義する。\par
      これを繰り返してゆくと一般に$\ds n ( \in \mathbb{N} )$に対して$\ds f_n$は弱くなる、つまり$\ds f_1 \geq f_2 \geq f_3 \geq \cdots$となってゆくので
      最終的に$\ds f$のすべての項を$\ds g_m$の和で表せる。($\ds m \in \mathbb{N}$)

      \begin{equation*}
        \begin{split}
          $\ds \eqref{eq:define-formula-of-f_1} \Longleftrightarrow f $ & $\ds = g_1(\sigma_1, \sigma_2, \cdots , \sigma_n) + f_1(x_1, x_2, \cdots , x_n)$\\
          &\begin{split}
            $\ds = g_1(\sigma_1, \sigma_2, \cdots , \sigma_n) + $ & $\ds g_2(\sigma_1, \sigma_2, \cdots , \sigma_n)$\\
            &$\ds + f_2(x_1, x_2, \cdots , x_n)$
          \end{split}\\
          &\vdots \\
          &\begin{split}
            $\ds = g_1(\sigma_1, \sigma_2, \cdots , $ & $\ds \sigma_n) + g_2(\sigma_1, \sigma_2, \cdots , \sigma_n)+$\\
            &$\ds \cdots + g_m(\sigma_1, \sigma_2, \cdots , \sigma_n)$
          \end{split}
        \end{split}
      \end{equation*}

      これにて一般に$\ds x_1, x_2, \cdots , x_n$に関する対称式$\ds f$が同対象に関する基本対称式のみで表せることが示されたが、
      基本対称式が一意に定まるか否かについての証明をしていないので、次にそれを示そう。\par
      $\ds f$を基本対称式より作られる対称式が複数あると仮定、つまり

      \begin{multline}
        \begin{split}
          $\ds \exists_{g, g'}$
          &\begin{cases}
            &\begin{cases}
              $\ds f(x_1, x_2, \cdots , x_n) $ & $\ds = g(\sigma_1, \sigma_2, \cdots , \sigma_n)$\\
              $\ds f(x_1, x_2, \cdots , x_n) $ & $\ds = g'(\sigma_1, \sigma_2, \cdots , \sigma_n)$
            \end{cases}\\
            &$\ds g(x_1, x_2, \cdots , x_n) \ne g'(x_1, x_2, \cdots , x_n)$\\
          \end{cases}
        \end{split}\\
        \label{eq:assum-multiple-equations-exist}
      \end{multline}

      を仮定する。\par
      ここで

      \begin{equation*}
        \begin{split}
          $\ds h(x_1, x_2, \cdots , x_n) \ldef $ & $\ds g(x_1, x_2, \cdots , x_n)$\\
          &$\ds - g'(x_1, x_2, \cdots , x_n)$
        \end{split}
      \end{equation*}
      
      と置くと仮定より

      \begin{numcases}
        $\ds h(x_1, x_2, \cdots , x_n) \ne 0$ & \label{eq:h-not-0} \\
        $\ds h(\sigma_1, \sigma_2, \cdots , \sigma_n) = 0$ & \label{eq:h-is-0}
      \end{numcases}

      である。\par
      一方で次のような$\ds \omega$の$\ds n$次方程式

      \begin{equation*}
        \begin{split}
          $\ds \omega^n - x_1 \omega^{n-1} + x_2 \omega^{n-2} - \cdots + (-1)^n x_n = 0$
        \end{split}
      \end{equation*}

      の解を$\ds \omega = \omega_1, \omega_2, \cdots , \omega_n$とすると解と係数の関係より

      \begin{cases}
        &$\ds \sigma_1(\omega_1, \omega_2, \cdots , \omega_n) = x_1$\\
        &$\ds \sigma_2(\omega_1, \omega_2, \cdots , \omega_n) = x_2$\\
        &\vdots \\
        &$\ds \sigma_n(\omega_1, \omega_2, \cdots , \omega_n) = x_n$\\
      \end{cases}
      \begin{equation*}
        \begin{split}
          $\ds (\because \eqref{eq:solution-and-coefficients})$
        \end{split}
      \end{equation*}

      これを\eqref{eq:h-is-0}に代入すると

      \begin{equation*}
        \begin{split}
          $\ds h(\sigma_1, \sigma_2, \cdots , \sigma_n) = h(x_1, x_2, \cdots , x_n)$
        \end{split}
      \end{equation*}

      となり、これは$\ds \eqref{eq:h-not-0} \land \eqref{eq:h-is-0}$に反する。\par
      よって仮定\eqref{eq:assum-multiple-equations-exist}が誤りで$\ds f$を基本対称式で表すその表し方はただ1通りである。
    \end{proof}
    
\end{document}